% !PRINTED 501
\weberpage{501}

Sind in dieser Darstellung Zähler und Nenner imprimitiv, so
kann man den Theiler herausnehmen und erhält eine Darstellung
in der Form
\weq{2}{A = a\,\frac{E_1(x,y,z,\ldots)}{E_2(x,y,z,\ldots)} = a\,\frac{E_1}{E_2} = aE,}
in der $a$ eine \gesperrt{positive}, ganze oder gebrochene rationale Zahl
ist, während $E_1, E_2$ primitive Functionen in $R$ sind. Hieraus
folgt:

1. \gesperrt{Man kann jedes rationale Functional durch
Multiplication mit einer primitiven Function in $R$ in eine
ganze Function in $R$ verwandeln, und mehrere rationale
Functionale lassen sich als Brüche darstellen, deren
gemeinsamer Nenner eine primitive Function ist.}

Die positive Zahl $a$ und der Quotient $E_1:E_2$ sind durch $A$
vollständig bestimmt. Denn setzen wir $a$ in die Form eines
rationalen Bruches $q_1:q_2$ und nehmen an, es sei
\[
\frac{q_1}{q_2}\,\frac{E_1}{E_2} = \frac{q_1'}{q_2'}\,\frac{E_1'}{E_2'},
\]
so folgt
\[
q_1 q_2' E_1 E_2' = q_2 q_1' E_2 E_1'.
\]
% [sic: Weber prints ``Functionen $R$'' without ``in''.]
Hier haben wir also zwei ganze Functionen $R$ mit ganzzahligen
Co\"efficienten, die einander gleich sind, und deren Theiler, da
$E_1 E_2'$ und $E_2 E_1'$ primitive Functionen sind, $q_1 q_2'$
oder $q_2 q_1'$ ist. Folglich ist $q_1 q_2' = q_2 q_1'$, und daher
auch
\[
\frac{q_1}{q_2} = \frac{q_1'}{q_2'}, \quad \frac{E_1}{E_2} = \frac{E_1'}{E_2'}.
\]

2. \gesperrt{Wir nennen die positive rationale Zahl $a$ den
absoluten Werth des Functionales $A$.}

In dem Falle, dass $A$ selbst eine Zahl ist, ist $a$ der absolute
Werth von $A$ in dem gewöhnlichen Sinne dieses Wortes, und $E$
ist $= +1$ oder $= -1$ zu setzen, je nachdem $A$ positiv oder
negativ ist. Es ist also $E$ nichts weiter als das Vorzeichen von
$A$. In der weitgehenden Verallgemeinerung dieser elementaren
Begriffe liegt das Befremdende, was unsere Definition dem ersten
Blick bietet. Sie wird sich aber in der Folge als durchaus
sachgemäss und nützlich erweisen.

Aus \S.\,135, 3. ergiebt sich der Satz:
%CONT>>
